\section{The finite ordinary alternating range}
\label{sec:ordinary-alternating-small}

This section treats ordinary representations in alternating degree
$8\le n\le14$.  It does not extrapolate a finite scan to larger $n$.

\begin{lemma}[HD-ALT-02a: the complete finite spectral input]
\label{HD-ALT-02a}
Let $8\le n\le14$.  Among nontrivial ordinary irreducible representations of
$\mathfrak A_n$, call its deleted natural representation, of degree
$n-1$, standard.  The following statements hold.
\begin{enumerate}
\item Every nonstandard irreducible representation has degree at
least fourteen.  Every nonscalar core element has $\nu\ge a/5$.
\item In every nonstandard representation, every nonscalar linear
normalizer element has $\nu\ge a/7$.
\item The only nonstandard degree at most fifteen is the degree-fourteen
$A_8$ character with Specht label $(4,4)$, up to transpose.  Its full
linear normalizer satisfies $\nu\ge a/4$.
\item Every nonreal core character has degree greater than fifteen.
Its linear normalizer is scalar times the core, and therefore satisfies
$\nu\ge a/5$.
\end{enumerate}
All scalar phases are allowed in these assertions.
\end{lemma}

\begin{proof}
We use the complete generic character tables
\texttt{CharacterTable("Alternating",n)} and
\texttt{CharacterTable("Symmetric",n)} from the published GAP
character-table implementation.  Their exact meaning and parameter
conventions are documented in \cite[\S4.2 and \S4.2-1]{SPOR:CTblLib}.
These are computed generic tables, not relabelled ATLAS tables.
The retained inputs contain every character, class size and order,
all prime power maps, the seven actual natural-subgroup fusions,
the complete partition parameters and the source strings.
The input SHA-256 is
\texttt{f406cf1ad89801369196f36422c4dc2608e5cf96c763e5fab53b8a2ecfa1939c}.

The exact GAP scan checks all nonlinear alternating rows and all
symmetric rows whose alternating restriction has norm one.  The
fourteen tables give respectively
\[
\begin{array}{c|rrrrrrr}
n&8&9&10&11&12&13&14\\\hline
A_n&13&17&23&30&42&54&71\\
S_n&18&26&38&52&72&96&130
\end{array}
\]
selected rows, totaling $682$.  For all $49\,650$ nonscalar entries
the exact cyclic Fourier formula computes every eigenspace,
checks nonnegative integral multiplicities and their degree sum.
The standalone integer-polynomial reconstruction in
\texttt{research/alternating\_small\_certificate.py} selects rows
again by pointwise equality of restricted characters, checks all
seven fusions and power maps, and reconstructs every spectrum.
It matches every minimum and every complete minimizing spectrum
from the GAP output.  This is a separate arithmetic implementation,
with the generic character-table algorithm retained as an explicit
mathematical input.

The complete checks give the first three numerical assertions
on all these rows.  In particular all nonstandard alternating
rows have $\nu/a\ge1/5$, all selected nonstandard symmetric
rows have $\nu/a\ge1/7$, and the degree-fourteen rows all have
$\nu/a\ge1/4$.  The only core degrees at most fifteen are
$n-1$ and the indicated extra fourteen at $n=8$.
The exact conjugation test on every character value shows that
all nonreal alternating rows have degree greater than fifteen;
every symmetric character value in these tables is real.

To pass from those tables to all normalizers, use
$\operatorname{Aut}(A_n)=S_n$ and Schur's lemma, as in
HD-ALT-01a.  An inner normalizer is scalar times a core element.
If an odd outer action preserves a core representation, an
intertwiner for a transposition has scalar square.  Rescaling
it makes its square one, giving an actual representation of
$A_n\rtimes C_2=S_n$ with irreducible core restriction.
The original normalizer matrix, after its inner adjustment,
is therefore scalar times an element of one of the selected
represented symmetric tables.  This includes arbitrary actual
odd elements, not only transpositions.

If the core character were nonreal, such a symmetric extension
would be impossible: all symmetric character values here are
real, so their restrictions are real.  Thus its normalizer
has only inner projective action, proving the fourth assertion.
Scalar multiplication preserves $\nu$ throughout.
\end{proof}

\begin{theorem}[HD-ALT-02b: the remaining small ordinary core]
\label{HD-ALT-02b}
Let a dominant nonsplit Weil sixfold finite-cover source have exact
imaginary quadratic original character field.  Suppose an inducing
block of degree at least five is primitive and factorwise tensor
indecomposable in the convention of D-BI-19.  If its forced core
is ordinary $A_n^t$ with $8\le n\le14$, then necessarily
\[
 n\in\{11,12,13,14\},\qquad t\ge2,
 \quad V|_{A_n^t}\simeq\operatorname{Std}_n^{\otimes t}.
                                                               \tag{ALTS.1}
\]
In particular no single-factor ordinary core in this finite range
can occur.  Larger inducing character fields and all compatible
factor permutations are included.
\end{theorem}

\begin{proof}
The core factors are permuted transitively by D-BI-19.  Consequently
their representations differ only by automorphisms of $A_n$.
Being standard and being self-dual are preserved by those
automorphisms, and all factor degrees are the same, say $a$.
Write $e=a^t$.

First suppose the factor characters are nonreal.  HD-ALT-02a gives
$a>15$ and a fifth moving-rank bound on every factor normalizer.
For a factor-preserving tensor operator that bound passes to the
total degree by fixing all other eigenvalues.  For a nontrivial
factor permutation the weighted cyclic-word argument gives the
stronger quarter bound.  Genuine scalar inertia has full moving
rank.  The actual intermediate-cover genus and branch budget
therefore gives
\[
                 12\lceil e/5\rceil\le2e+6,
                 \qquad e\le15,
\]
contrary to $e\ge a>15$.

The factor characters are thus real, hence the core representation
is self-dual.  Its invariant nondegenerate bilinear form is unique
up to scalar, so the full inducing group is essentially self-dual.
Suppose its factors are nonstandard.  Their degree is at least
fourteen, so the primitive ESD bound D-BI-09 applies.  At least
two genuine inertias move at least half the total degree; the
other ten obey the seventh bound of HD-ALT-02a, including tensor
permutations because $1/4>1/7$.  Hence
\[
                   e+10e/7\le2e+6,
                   \qquad e\le14.                         \tag{ALTS.2}
\]
The finite degree list forces $t=1$, $n=8$ and $a=e=14$.
Its stronger quarter normalizer bound instead gives
$12\lceil14/4\rceil=48>34=2e+6$, a contradiction.

It remains that all factors are standard.  If $t=1$, the complete
scalar deleted-standard theorem D-BI-08 excludes this nonsplit
source.  If $t\ge2$ then $e\ge(n-1)^2\ge49$, so D-BI-09 applies.
Every nonscalar factor normalizer has $\nu\ge1=a/(n-1)$,
and factor permutations satisfy the stronger quarter bound.
The same two-half-inertia budget gives
\[
             e+\frac{10e}{n-1}\le2e+6.
\]
For $n=8,9,10$ it implies respectively $e\le14,24,54$,
whereas two factors already have degree at least $49,64,81$.
These cases are impossible.  Precisely the stated necessary
conditions (ALTS.1) remain.  No assertion of their existence or
dominance is made.
\end{proof}

\paragraph{Limits.}
This theorem concerns the ordinary core.  Spin cores in $n\ge8$
were excluded separately by HD-SPIN-02.  Exceptional covers with
$n<8$ and ordinary alternating degree $n\ge15$ retain their own
separate restrictions.  In (ALTS.1) the scalar deleted-standard
single-factor theorem is not applied to a tensor power.
